Nesta seção, é realizada uma revisão dos principais trabalhos da literatura que investigam a aplicação de métodos computacionais na análise de imagens histopatológicas para o diagnóstico do câncer de próstata, evidenciando as abordagens adotadas, os conjuntos de dados utilizados, os resultados obtidos e as lacunas identificadas.

No trabalho de \cite{Xiang2023Prostate}, os autores propuseram a ferramenta GCN-MIL, que utiliza uma Rede Neural de Grafos para agregar características espaciais de \textit{patches} em lâminas de WSI. O modelo incorpora a estratégia de \textit{Robust Training}, que filtra iterativamente amostras de alta incerteza para mitigar o impacto de rótulos ruidosos e inconsistentes, alcançando um Kappa quadrático de 93,10\% no conjunto interno do desafio PANDA. Entretanto, a abordagem apresenta limitações metodológicas relevantes. O emprego da função de perda \textit{Cross-Entropy} tradicional, que trata as graduações ISUP como classes nominais independentes, ignora a hierarquia biológica inerente à progressão do câncer de próstata, o que pode comprometer a coerência ordinal das predições. Ademais, o uso da ResNet50 como extrator de características, uma arquitetura clássica de RNC, pode restringir tanto a capacidade de representação morfológica quanto a eficiência computacional em comparação com modelos mais modernos.

Em \cite{kosoko2024xai_prostate}, os autores avaliaram diferentes arquiteturas de RNCs, combinadas à técnica de \textit{Grad-CAM}, para a interpretação das decisões do modelo. O conjunto de dados utilizado foi o SICAPv2, composto por \textit{patches} extraídos de WSIs coradas em H\&E, com anotação em nível de pixel, totalizando 4417 amostras não cancerosas: 1635 do grupo GG3, 3622 do GG4 e 665 do GG5. Entre as arquiteturas avaliadas, o modelo VGG19 apresentou o melhor desempenho, alcançando AUC-ROC de 85,0\% e métricas médias de acurácia, precisão e \textit{recall} em torno de 75,0\%. Entretanto, o estudo apresenta lacunas significativas. A ausência de validação clínica externa compromete a avaliação da capacidade de generalização do modelo em cenários interinstitucionais. Ademais, o emprego exclusivo de arquiteturas clássicas de RNCs restringe a capacidade de representação morfológica e o desempenho em conjuntos de dados mais heterogêneos.

No estudo de \cite{afifi2024prostate}, a detecção de câncer de próstata foi formulada como um problema de classificação binária, no qual foram avaliadas as arquiteturas InceptionV3, ResNet50 e InceptionResNetV2 em conjunto com as técnicas de \textit{Transfer Learning} e \textit{Fine-Tuning}, sobre imagens WSI provenientes do \textit{Radboud University Medical Center} e do \textit{Karolinska Institute}. O modelo InceptionResNetV2 apresentou o melhor desempenho, atingindo acurácia de 91,80\%. Entretanto, a formulação binária ignora a graduação histológica da doença, o que limita a utilidade clínica do modelo, uma vez que a distinção entre graus de malignidade é determinante para o planejamento terapêutico. Ademais, o uso exclusivo da acurácia é insuficiente para uma avaliação robusta, especialmente em contextos de desequilíbrio entre as classes.

Em \cite{alici2024efficient}, os autores incorporaram mecanismos de atenção à arquitetura EfficientNet-B4 para o diagnóstico e a graduação do câncer de próstata, utilizando o conjunto de dados público DiagSet. A abordagem contempla três níveis de classificação: binária (presença ou ausência de câncer), quatro classes (tecido normal, estroma, alterações benignas e câncer) e oito classes (tecido normal, estroma, artefatos, inflamação, hiperplasia benigna, Gleason 3, 4 e 5). Os resultados apresentam acurácias de 96,18\%, 94,86\%, 93,32\%, respectivamente. Entretanto, o uso exclusivo de fragmentos de WSI limita a representatividade do cenário clínico real, e o reporte restrito à acurácia é insuficiente para uma avaliação robusta diante do desequilíbrio entre as classes. Ademais, a aplicação da classificação binária limita a utilidade clínica do modelo. 

De maneira geral, os trabalhos revisados apresentam algumas limitações em comum. No que diz respeito à avaliação de desempenho, os estudos de \cite{afifi2024prostate, alici2024efficient} reportam apenas a acurácia, uma métrica insuficiente para uma avaliação robusta em cenários de desbalanceamento entre as classes. O emprego de RNCs clássicas, como ResNet50 e VGG19, como extratores de características \cite{Xiang2023Prostate, kosoko2024xai_prostate} pode restringir a capacidade de representação morfológica em comparação com arquiteturas recentes. Ademais, estudos que adotam classificação binária \cite{afifi2024prostate, alici2024efficient} ignoram a graduação histológica da doença, o que limita a utilidade clínica dos modelos, já que a distinção entre graus de malignidade é determinante para o planejamento terapêutico.

Com o intuito de superar as limitações identificadas, o presente trabalho propõe uma metodologia estruturada que contempla o tratamento de rótulos ruidosos, a utilização de arquiteturas de \textit{Deep Learning} e de estratégias de \textit{Ensemble} para potencializar a robustez das inferências. Além disso, a natureza ordinal da progressão biológica da doença é incorporada à modelagem por meio de uma função de perda sensível à hierarquia entre os grupos de grau ISUP, o que promove maior coerência e utilidade clínica das predições. Por fim, os modelos são avaliados por meio de um conjunto de métricas capaz de refletir, de forma mais robusta, seu desempenho.